\section{Martingale representation and optimization consequences}
\label{sec:consequences}
Full support gives the gradient martingale exactly the second-moment
profile of Brownian motion: $\E u_\gamma(t,X_t)^2=t$. Its correlation with
the driving Brownian motion recovers the ground-state value. We first
prove this identity and then approximate it using a single
$\{-1,1\}$-valued function of finitely many independent Gaussians.
The gradient-martingale and normalized Euler approach is developed in
Montanari~\cite[Sections~2--3]{Montanari2019} and El Alaoui, Montanari,
and Sellke~\cite[Sections~5--6 of the cited preprint]{EAMS2021}.
Here Theorem~\ref{thm:full-support} supplies the required zero-temperature
second-moment normalization.

\begin{corollary}[Martingale representation of the value]
\label{cor:value}
Let $\gamma$ be a minimizer, let $X$ be its diffusion, and write
$u=\Phi_{\gamma,x}$ and $a=\Phi_{\gamma,xx}$. There is a bounded
martingale terminal value $U_1$ such that $u(t,X_t)\to U_1$ almost surely
and in $L^2$ as $t\uparrow1$, with $U_1^2=1$ almost surely. Moreover,
\begin{equation}\label{eq:value-identities}
 \begin{aligned}
 P_*&=\E|X_1|-\int_0^1t\gamma(t)\,\dd t\\
    &=\E[U_1 B_1]
     =\int_0^1\E a(t,X_t)\,\dd t.
 \end{aligned}
\end{equation}
The last integral is finite.
\end{corollary}

\begin{proof}
Lemma~\ref{lem:diffusion} gives the bounded martingale
$u(t,X_t)=\int_0^t a(s,X_s)\,\dd B_s$. By
Theorem~\ref{thm:full-support}, its second moment is $t$ and
$\E a(t,X_t)^2=1$. Martingale convergence therefore gives $U_1$ with
$|U_1|\leq1$ and $\E U_1^2=1$, hence $U_1^2=1$ almost surely.
The diffusion converges to $X_1$ almost surely and in $L^2$, because its
drift has deterministic integrable bound $\gamma$.

The uniform terminal bound~\eqref{eq:terminal-uniform} and convexity
show that $u(t,x_t)\to\operatorname{sgn}(x)$ whenever $x_t\to x\ne0$.
Indeed, for a fixed small $h>0$, the derivative of a convex function at
$x_t$ lies between its left and right difference quotients of step $h$;
uniform convergence to $|\cdot|$ makes both limits equal to the sign
when the interval stays away from zero. Consequently
$U_1X_1=|X_1|$ almost surely, with no need to exclude $X_1=0$.

It\^o's formula along the optimal diffusion, justified on compact strips
as in Lemma~\ref{lem:diffusion}, gives
\[
 \E\Phi_\gamma(t,X_t)
  =\Phi_\gamma(0,0)+\frac12\int_0^t\gamma(s)
                                  \E u(s,X_s)^2\,\dd s.
\]
Let $t\uparrow1$. The terminal bound and $L^1$ convergence of $X_t$
identify the left side with $\E|X_1|$. Substitute
$\E u(s,X_s)^2=s$ and the definition of $\mathcal P$ to obtain the
first line of~\eqref{eq:value-identities}.
Since $X_1=B_1+\int_0^1\gamma(s)u(s,X_s)\,\dd s$, martingality
and Fubini's theorem yield
\[
 \E[U_1X_1]=\E[U_1B_1]
       +\int_0^1\gamma(s)\E[u(s,X_s)^2]\,\dd s.
\]
All terms are integrable because $|u|,|U_1|\leq1$ and $\gamma\in L^1$.
Finally the $L^2$ stochastic-integral identity at time one implies
$\E[U_1B_1]=\int_0^1\E a(s,X_s)\,\dd s$.
The integral is absolutely bounded by one using
$\E|a(s,X_s)|\leq(\E a(s,X_s)^2)^{1/2}=1$.
\end{proof}

The last expression in~\eqref{eq:value-identities} is the expected
covariation of the gradient martingale with its driving Brownian motion.
To approximate it, replace the diffusion by an Euler scheme on a fixed
strip $[0,T]$, where $T<1$. The continuum identity $\E a(t,X_t)^2=1$
makes each martingale increment have variance equal to its time step.
The Euler scheme has only an approximate version of this normalization,
so we normalize its Gaussian-weighted directions to unit variance.

One additional independent Gaussian then rounds the Euler terminal
gradient to a spin while preserving its conditional mean. The coefficients
$A_j$ below measure this spin's correlations with the normalized
martingale directions; the coefficients $B_i$ measure its correlations
with the original Gaussian coordinates. Their adjacent-index product
converges to the same covariation integral. These are scalar Gaussian
coefficients, not an implementation on the finite-spin disorder.

On this fixed strip, $u,a$ and all their spatial derivatives are bounded;
the Gaussian increments remain unbounded. We first refine the mesh and
only afterwards let $T$ approach one. The mesh count $N$ is distinct
from the spin-system size $n$ in the introduction.

\begin{proposition}[Finite Gaussian coefficients]\label{prop:gaussian}
Fix $0<T<1$, a minimizer $\gamma$, and $u,a$ as above. For an integer
$N\geq1$ put $\Delta=T/N$ and $t_j=j\Delta$. On independent standard
normal variables $Z_1,\ldots,Z_{N+1}$, define
\begin{align}
 Y_0&=0,\qquad
 Y_j=Y_{j-1}+\sqrt\Delta Z_j
      +\Delta\gamma(t_{j-1})u(t_{j-1},Y_{j-1}),
       &&1\leq j\leq N,\label{eq:euler}\\
 c_j&=\bigl(\E a(t_{j-1},Y_{j-1})^2\bigr)^{-1/2}>0,
 \qquad g_j=c_j a(t_{j-1},Y_{j-1})Z_j.
       \label{eq:normalized-increment}
\end{align}
For all sufficiently large $N$ these constants are finite and positive.
For these mesh counts, let $F_{\rm G}$ denote the standard normal
distribution function and set
\[
 \begin{gathered}
 F_N=\operatorname{sgn}\bigl(u(T,Y_N)+2F_{\rm G}(Z_{N+1})-1\bigr),\\
 A_j=\E[F_Ng_j],\qquad B_i=\E[F_NZ_i],\qquad A_0=0.
 \end{gathered}
\]
Here $1\leq j\leq N$ and $1\leq i\leq N+1$ in the definitions of
$A_j$ and $B_i$.
The value of the sign at zero is immaterial. Then
\begin{equation}\label{eq:finite-value}
 \lim_{N\to\infty}\sum_{j=0}^N A_jB_{j+1}
       =\int_0^T\E a(t,X_t)\,\dd t.
\end{equation}
Letting $T\uparrow1$ on the right gives $P_*$.
\end{proposition}

\begin{proof}
\emph{Euler convergence and normalization.}
Use Brownian increments $\sqrt\Delta Z_j=B_{t_j}-B_{t_{j-1}}$ to
couple~\eqref{eq:euler} to $X$ up to time $T$.
Lemma~\ref{lem:regularity}, on a slightly larger strip, bounds all
spatial derivatives of $u,a$. Their almost-everywhere time derivatives
given by the PDE are bounded there, so $u,a$ are uniformly Lipschitz
in time and space. The drift $\gamma u$ is spatially Lipschitz with
a common bound. For its time dependence, it is enough to use that
$\gamma$ is bounded and monotone.
If $\pi_N(s)=t_{j-1}$ on $[t_{j-1},t_j)$, then
\[
 \int_0^T|\gamma(s)-\gamma(\pi_N(s))|\,\dd s
       \leq\Delta\bigl(\gamma(T)-\gamma(0)\bigr).
\]
Also $\E|X_s-X_{\pi_N(s)}|^2\leq C\Delta$.
Subtract the two integral drift equations at the grid points, use these
bounds and the Lipschitz bounds, and apply discrete Gronwall. This gives
\begin{equation}\label{eq:euler-convergence}
 \E\max_{j\leq N}|Y_j-X_{t_j}|^2\longrightarrow0.
\end{equation}
More explicitly, the $L^2$ norm of the sum of local drift errors is at
most $C\sqrt\Delta+C\Delta$; the remaining errors satisfy
$e_j\leq C\Delta\sum_{i<j}e_i+o(1)$, which proves the claim.

Uniform convergence in~\eqref{eq:euler-convergence}, boundedness and
Lipschitz continuity of $a$, and $\E a(t,X_t)^2=1$ imply
$\max_{j\leq N}|c_j-1|\to0$.
The $g_j$ form an orthonormal family in $L^2$: they have unit second
moment by definition and zero conditional mean given the previous
increments. It\^o isometry,~\eqref{eq:euler-convergence}, and the
same continuity estimates give
\[
 M_N:=\sum_{j=1}^N\sqrt\Delta\,a(t_{j-1},Y_{j-1})Z_j
      =u(T,Y_N)+o_{L^2}(1).
\]
\emph{Rounding and the two coefficient vectors.}
The independent random variable $F_{\rm G}(Z_{N+1})$ is uniform on
$(0,1)$. For any $v\in[-1,1]$, the sign of $v+2U-1$, with $U$
uniform on $(0,1)$, has mean $v$. Hence
$\E[F_N\mid Z_1,\ldots,Z_N]=u(T,Y_N)$ exactly.
This preserves every correlation with a function of the first $N$
Gaussians. Bessel's inequality applied to $u(T,Y_N)-M_N$ against the two
orthonormal families $(g_j)_{j=1}^N$ and $(Z_i)_{i=1}^N$ yields
\begin{align*}
 (A_j)_{j=1}^N&=(\sqrt\Delta/c_j)_{j=1}^N+o_{\ell^2}(1),\\
 (B_i)_{i=1}^N&=(\sqrt\Delta\,\E a(t_{i-1},Y_{i-1}))_{i=1}^N
                                      +o_{\ell^2}(1).
\end{align*}
For the second equality, terms of $M_N$ before $i$ are independent of
$Z_i$, and those after $i$ have zero conditional mean; the $i$th term
contributes the stated expectation. The first equality follows from
orthogonality of the $g_j$ and $M_N=\sum_j\sqrt\Delta\,g_j/c_j$.

\emph{The shifted product.}
The first coefficient vector is asymptotically the constant vector
with entries $\sqrt\Delta$, because $c_j\to1$ uniformly. The second
has entries $\sqrt\Delta$ times the expected Hessian at the preceding
mesh point. Pairing $A_j$ with $B_{j+1}$ therefore selects the expected
Hessian at $t_j$, which is a Riemann-sum value for the target integral.
To justify the accumulated errors, the vectors in both displays have
bounded $\ell^2$ norm. Cauchy--Schwarz shows that their errors contribute
$o(1)$ to the shifted scalar product with $1\leq j<N$. The first display also gives
$A_N=\sqrt\Delta/c_N+o(1)=o(1)$. The remaining term therefore obeys
$|A_NB_{N+1}|\leq|A_N|=o(1)$, since $|B_{N+1}|\leq1$.
Thus the sum in~\eqref{eq:finite-value} differs by $o(1)$ from
\[
 \Delta\sum_{j=1}^{N-1}\frac{\E a(t_j,Y_j)}{c_j}.
\]
This converges to the stated integral by~\eqref{eq:euler-convergence},
uniform convergence of $c_j$ to one, and continuity of
$t\mapsto\E a(t,X_t)$. Corollary~\ref{cor:value} gives its limit as
$T\uparrow1$.
\end{proof}
